\section{Related Work}
\label{sec:related_work}

Early GPU AES studies established the feasibility of exploiting block-level parallelism
\cite{manavski2007cuda,harrison2007aes}. Subsequent work examined thread granularity, memory
placement and overlapping computation with transfers~\cite{iwai2012acceleration}. Lookup-table,
shared-memory and bit-sliced implementations further demonstrate how representation and memory
access affect GPU throughput~\cite{wang2019gpu,tezcan2021optimization,bitsliced2024gpu}. These
studies motivate GPU acceleration, but their rates are specific to the implementation,
hardware and measured interval. Kernel optimization results cannot be treated as interchangeable
with transfer-inclusive measurements.

The CPU reference is equally important. AES-NI accelerates AES transformations through dedicated
instructions~\cite{intel2012aesni}. More generally, Lee et al.~\cite{lee2010debunking} show why
CPU--GPU comparisons depend on the implementations selected as references. A gain over an
unaccelerated CPU path therefore does not establish a gain over AES-NI. The present study uses
the latter as its CPU denominator, while recognizing that a selected library path is not an
upper bound on host performance.

Data location and movement are established concerns in heterogeneous benchmarking.
Gregg and Hazelwood~\cite{gregg2011data} emphasize their role in interpreting CPU--GPU results;
AES-specific research has also considered transfer overlap~\cite{iwai2012acceleration}.
CUDA and CuPy guidance additionally require suitable synchronization when timing asynchronous
device operations~\cite{nvidiaBestPractices,cupyPerformance}. The distinction between resident
execution and a broader processing path is therefore not introduced here as a novel concept.

Reliable reporting also requires a suitable replication structure. Mytkowicz et al.
\cite{mytkowicz2009wrong} demonstrate how experimental conditions can bias measurements.
Kalibera and Jones~\cite{kalibera2013rigorous} address variability across replication levels,
while Hoefler and Belli~\cite{hoefler2015scientific} discuss defensible reporting in parallel
computing. These principles motivate paired session estimates and explicit limits on inference.

The gap addressed here is an integrated, auditable assessment of when measured GPU advantages
survive both an AES-NI reference and a transfer-inclusive boundary across the selected systems.
The study combines established benchmarking principles with cryptographic validation and
archived evidence, reporting consistent gains alongside reversals and variable outcomes. It
complements implementation-optimization studies rather than attempting to rank their kernels
or reproduce their performance records.
